% Fichier complété par tools/complete_missing_corrections.py.
% Source : CIN/CIN-01-Parametrage/07_RR3D/07_RR3D.tex
% Statut : rédigé (1 question(s)).
\begin{corrigebox}[Corrigé rédigé]
\CorrigeQuestion{1}
On choisit $\theta(t)$ pour la rotation 1/0 autour de $\vec k_0$, puis
$\varphi(t)$ pour la rotation 2/1 autour de $\vec i_1$. On a
\[\vec\Omega_{2/0}=\dot\theta\vec k_0+\dot\varphi\vec i_1,\qquad
  \overrightarrow{AB}=R\vec i_1,\qquad
  \overrightarrow{BC}=\ell\vec i_1+r\vec j_2.\]
Les changements de base sont une rotation d'angle $\theta$ autour de
$\vec k_0$, puis une rotation d'angle $\varphi$ autour de $\vec i_1$.
\end{corrigebox}
